% 本文件由 LaTeX 排版 Agent 根据有效 translation.json 自动生成。
% author=Clement of Rome; work_id=1011; title=圣克肋孟后书

\chapter{圣克肋孟后书}

% structure: p0001 source=h1 target=omitted reason=page-title-already-rendered-by-parent

% structure: p0002 source=h2 target=section reason=relative-heading-rank; base=h2

\section{第一章．我们当对基督有崇高思想}

弟兄们，你们应当把耶稣基督看作天主来思念——作为审判活人与死者的。我们不应该轻看我们的救恩；因为如果我们对祂估计很低，我们也就只能希望从祂那里得到少许。我们中那些漫不经心地听这些事，仿佛它们不重要的人，是在犯罪，不知道我们是从哪里被召来的，被谁所召，召到什么境地，也不知道耶稣基督为我们的缘故忍受了多少苦难。那么，我们该用什么来回报祂，或者结出什么与祂所赐给我们相配的果实呢？确实，我们欠祂的恩宠是何等巨大！祂仁慈地赐给了我们光明；作为父亲，祂称我们为儿子；在我们将要灭亡时，祂拯救了我们。那么，我们该给祂什么赞美，或用什么来回报我们所领受的恩赐呢？我们缺乏理智，崇拜石头、木头、金子、银子、铜器，以及人手的作品；我们的整个生命无非就是死亡。我们被盲目所困，眼前一片黑暗，却得以看见，并因祂的旨意脱去了笼罩我们的云彩。因为祂怜悯了我们，仁慈地拯救了我们，看我们纠缠在许多错误中，并面临毁灭，而且我们除了从祂那里得到救恩以外，毫无希望。因为祂在我们不存在时召叫了我们，并愿意我们从虚无中达到真实的存在。\par

% structure: p0004 source=h2 target=section reason=relative-heading-rank; base=h2

\section{第二章．教会，从前不生育，如今多结果实}

\BibleQuote{不生育的，你要喜乐；不生产的，你要欢呼；因为被遗弃的比有丈夫的儿女更多。}祂说：\BibleQuote{不生育的，你要喜乐}这句话指的是我们，因为我们的教会过去不生育，直到有儿女赐给她。但当祂说：\BibleQuote{不生产的，你要欢呼}意思是，我们应当真诚地向天主献上祈祷，不应像产妇那样显出软弱。而当祂说：\BibleQuote{因为被遗弃的比有丈夫的儿女更多}[祂的意思是]我们的人民似乎被天主弃绝，但现在通过信仰，比那些被认为拥有天主的人更多。另一处经上说：\BibleQuote{我来不是召义人，而是召罪人。}这意味着那些走向灭亡的人必须被拯救。因为建立稳固的根基不算什么，扶持那将要倒塌的才是伟大而可赞叹的事。基督也如此愿意拯救那些将要灭亡的\BibleRef{玛 18:11}，祂来临并召叫我们这些奔向毁灭的人，拯救了许多人。\par

% structure: p0006 source=h2 target=section reason=relative-heading-rank; base=h2

\section{第三章．宣认基督的本分}

既然祂向我们显出了如此大的慈悲，尤其是在这一点上：我们这些活着的人不应向死了的神祇献祭或朝拜它们，而应通过祂获得对真正圣父的认识。那么，我们如何表明我们真的认识祂呢？不是通过否认我们藉以获得这认识的那位吗？因为祂亲自宣称：\BibleQuote{凡在人前承认我的，我在我天父前也必承认他。}\BibleRef{玛 10:32}那么，如果我们承认我们得救所依靠的那位，这就是我们的赏报。但我们该如何承认祂呢？通过做祂所说的事，不违反祂的诫命，并不仅用嘴唇，而要用全心和全意来尊敬祂。\BibleRef{玛 22:37}因为祂在依撒意亚书中说：\BibleQuote{这些人用嘴唇尊敬我，他们的心却远离我。}\BibleRef{依 29:13}\par

% structure: p0008 source=h2 target=section reason=relative-heading-rank; base=h2

\section{第四章．对基督的真实宣认}

让我们不要仅仅称祂为主，因为那样不能救我们。因为祂说：\BibleQuote{不是每一个向我说主啊，主啊的人，都能得救，而是那行义的人。}所以，弟兄们，让我们通过我们的行为来宣认祂：彼此相爱、不犯奸淫、不彼此毁谤、不怀嫉妒，而是节制、怜悯、良善。我们也应当彼此同情，不贪婪。通过这样的行为让我们宣认祂，而不是通过相反的行为。我们不应该怕人，而应怕天主。因此，如果我们做了这样的坏事，主曾说：\Quote{即使你们聚集在我的怀中，如果你们不遵守我的诫命，我也要把你们抛弃，并对你们说：离开我吧，我不认识你们从哪里来，你们这些作恶的人。}\par

% structure: p0010 source=h2 target=section reason=relative-heading-rank; base=h2

\section{第五章．应当轻视这个世界}

因此，弟兄们，让我们心甘情愿地离开今世的寄居，做那召叫我们的天主的旨意，不要害怕离开这个世界。因为主说：\BibleQuote{你们要像羊进入狼群。}\BibleRef{玛 10:16}伯多禄回答祂说：\Quote{那么，如果狼撕碎羊羔，怎么办呢？}耶稣对伯多禄说：\Quote{羊羔死后不必再怕狼；同样，不要怕那些能杀死你，却不能再做什么的人；而要怕那在你死后，有权柄把灵魂和身体都投入地狱之火的那一位。}弟兄们，请思考：在这世界肉体中的寄居是短暂而暂时的，但基督的许诺是伟大而奇妙的，即将来国度的安息和永生。那么，我们该通过什么行为才能获得这些呢？岂不就是过着圣洁正义的生活，把世俗之事看作不属于我们，不把欲望寄托其上吗？因为如果我们贪求占有它们，就会偏离正义之路。\par

% structure: p0012 source=h2 target=section reason=relative-heading-rank; base=h2

\section{第六章．今世与来世彼此为敌}

主宣告说：\Quote{没有仆人能事奉两个主人。}\BibleRef{玛 6:24}因此，如果我们企图同时事奉天主和\OriginalTerm{玛门}{mammon}，这对我们毫无益处。\Quote{人纵然赚得了全世界，却赔上了自己的灵魂，为他有什么益处？}\BibleRef{玛 16:26}今世与来世是两个仇敌。一个催逼人犯奸淫和败坏、贪婪和欺诈；另一个却与这些事告别。因此，我们不能同时作两者的朋友；我们必须舍弃一个，确保另一个。我们认为，憎恨眼前的事物是更好的，因为它们是短暂的、易逝的、可朽坏的；而爱慕将来的事物，因为它们美善而不可朽坏。如果我们遵行基督的旨意，我们必得安息；否则，如果我们违背祂的诫命，什么也不能拯救我们脱离永罚。因为圣经在厄则克耳中也这样说：\Quote{倘若诺厄、约伯和达尼尔兴起，他们也不能救自己的儿女脱离被掳。}\BibleRef{则 14:14}既然如此，这些如此义德的人尚且不能以自己的义德拯救他们的儿女，我们又如何能希望进入天主的王宫，除非我们保持我们的圣洗圣事圣洁无瑕？或者，谁将成为我们的辩护者，除非我们发现拥有圣善和义德的行为？\par

% structure: p0014 source=h2 target=section reason=relative-heading-rank; base=h2

\section{第七章 我们必须努力奋斗，以获得冠冕}

所以，我的弟兄们，让我们竭尽全力搏斗，知道竞赛就在我们眼前，有许多人远涉重洋去争取一个可朽坏的奖赏；然而并非所有人都得冠冕，只有那些刻苦努力、光荣奋战的人。因此，让我们也这样奋斗，好使我们都得冠冕。让我们奔跑正直的赛道，就是那不朽坏的赛跑；让我们大批地出发，奋力争取冠冕。即使我们不能全部得到冠冕，至少也要接近它。我们必须记住，在可朽坏的竞赛中，如果有人被发现行为不公，就会被带走、鞭打，并被逐出赛场。那么你们想，如果有人在不可朽坏的竞赛中作了不体面的事，他要承受什么？因为关于那些不保存印记的人，经上说：\Quote{他们的虫子不死，他们的火也不灭，他们将成为众人憎恶的对象。}\BibleRef{依 66:24}\par

% structure: p0016 source=h2 target=section reason=relative-heading-rank; base=h2

\section{第八章 我们在世上时悔改的必要性}

因此，只要我们还在世上，让我们实行悔改，因为我们在工匠手中如同泥土。因为正如陶工，如果他制造一个器皿，而在手中变得扭曲或破裂，便重新塑造它；但如果他已经把它投入火窑，就不能再为它找到任何帮助。同样，让我们也这样，当我们还在这个世界时，全心全意地悔改我们在肉身中所行的恶事，好让我们还有悔改机会时，被主拯救。因为离开世界之后，我们不再有告解或悔改的能力。所以，弟兄们，通过遵行父的旨意，保持肉身圣洁，遵守主的诫命，我们将获得永生。因为主在福音中说：\Quote{如果你没有保管好那微小的，谁会把大事交给你呢？因为我告诉你们：那在最小的事上忠信的，在大事上也忠信。}\BibleRef{路 16:10}这就是祂的意思：\Quote{保持肉身圣洁，印记无瑕，好使你们获得永生。}\par

% structure: p0018 source=h2 target=section reason=relative-heading-rank; base=h2

\section{第九章 我们将以肉身受审判}

你们中不可有人说：这肉身不会被审判，也不会复活。要想一想，你们是在什么状态中被拯救的，你们是在什么状态中看见光明的，如果不是在这肉身中呢？因此，我们必须保存肉身如同天主的殿。因为正如你们在肉身中被召叫，你们也要在肉身中来受审判。正如基督主，祂拯救了我们，祂先是灵，然后成了肉身，并这样召叫了我们，同样，我们也要在这肉身中领受赏报。因此，让我们彼此相爱，好使我们都能达到天主的国。当我们还有机会被医治时，让我们把自己交托给医治我们的天主，并回报祂。什么样的回报呢？出自真诚之心的悔改，因为祂预知一切，并知晓我们心中的事。因此，让我们赞美祂，不只用口，也用心灵，好使祂接纳我们作儿子。因为主曾说：\Quote{那些遵行我父旨意的人，就是我的弟兄。}\BibleRef{玛 12:50}\par

% structure: p0020 source=h2 target=section reason=relative-heading-rank; base=h2

\section{第十章 要弃绝恶习，追随美德}

所以，我的弟兄们，让我们遵行那召叫我们的父的旨意，好使我们生活；让我们热切追随美德，但要弃绝每一个引我们犯罪的邪恶倾向；要逃避不虔敬，免得祸患临到我们。因为如果我们勤于行善，平安必随我们而来。因此，那些受人的恐惧影响，宁愿选择眼前的享乐而不顾将来应许的人，不能找到平安。因为他们不知道眼前的享乐带来多大的痛苦，也不知道将来的应许包含多大的幸福。而且，如果只是他们自己这样做，那还可以容忍；但现在他们坚持用有害的教义污染无辜的灵魂，不知道他们和那些听他们的人都将受到双重的惩罚。\par

% structure: p0022 source=h2 target=section reason=relative-heading-rank; base=h2

\section{第十一章 我们应当事奉天主，信赖祂的应许}

所以，让我们以纯洁的心事奉天主，我们便成为义人；但如果我们不信天主的许诺而不事奉祂，我们就将是不幸的。因为先知的话也宣告说：\Quote{三心二意、内心怀疑的人是有祸的，他们说：\InnerQuote{这一切我们在祖先的时代就已经听说了；但我们日日等待，却没有看见任何一样成就。}愚昧的人哪！把你们自己比作一棵树吧，就拿葡萄树来说。它先是落叶，然后发芽，接着是酸葡萄，最后是成熟饱满的果实。同样，我的子民也经历了扰乱和苦难，但以后他们必将获得福乐。}因此，我的弟兄们，我们不要三心二意，而要怀抱希望、坚忍到底，好使我们也能获得赏报。因为应许按各人的行为赐予赏报的那一位是信实的。所以，如果我们在天主面前行正义，我们必将进入祂的国，并领受那\Quote{眼所未见，耳所未闻，人心也未曾想到的}应许。\BibleRef{格前 2:9}\par

% structure: p0024 source=h2 target=section reason=relative-heading-rank; base=h2

\section{第十二章  我们当不断期待天主的国度}

因此，我们当在爱德和正义中，时时期待天主的国度，因为我们不知道天主显现的日子。因为主曾被人问及祂的国何时来到，祂回答说：\Quote{当二成为一，外在如同内在，男与女既非男亦非女时。}所谓二成为一，就是当我们彼此说真话，两个身体内毫无虚假地只有一个灵魂。而\Quote{外在如同内在}的意思是指：祂称灵魂为\Quote{内在}，身体为\Quote{外在}。所以，正如你的身体是肉眼可见的，也当让你的灵魂藉善行显明出来。至于\Quote{男与女，既非男亦非女}，祂这样说，意思是弟兄看见姊妹时，不可对她存有关于女性之念；姊妹看见弟兄时，也不可对他存有关于男性之念。祂说：\Quote{你们若做这些事，我父的国就要来临。}\BibleRef{格前 7:29}\par

% structure: p0026 source=h2 target=section reason=relative-heading-rank; base=h2

\section{第十三章  天主的圣名不应受亵渎}

所以，弟兄们，如今我们终当悔改，清醒地转向善事；因为我们充满了极多的愚妄和邪恶。让我们洗除自己先前的罪，从心里悔改而得救；我们不要作讨人喜欢的人，也不只愿彼此取悦，也要为了正义的缘故取悦外人，免得圣名因我们而受亵渎。因为主说：\Quote{我的名字在万民中不断地受亵渎}，又说：\Quote{因此我的名字受亵渎：因何受亵渎？因你们不做我所愿意的事。}\BibleRef{依 52:5} 因为外邦人从我们口中听到天主的圣言，对其卓越和珍贵感到惊奇；但随后得知我们的行为与所说的话不相称——他们便趁机转向亵渎，说那只是神话和迷惑。因为每当他们听我们说，天主说：\Quote{你们若爱那爱你们的人，便没有功劳；但你们若爱你们的仇敌和恨你们的人，便有功劳}——他们听到这话，便惊叹这善举的无比；但见到我们不只不爱那恨我们的人，甚至不爱那爱我们的人，他们就嘲笑我们，圣名便受了亵渎。\par

% structure: p0028 source=h2 target=section reason=relative-heading-rank; base=h2

\section{第十四章  属神的教会}

所以，弟兄们，若我们承行父天主的旨意，我们便成为那最初的、属神的教会——即在太阳和月亮之前受造的教会的肢体；但若我们不行主的旨意，我们便应验了圣经所说的：\Quote{我的家成了贼窝。}\BibleRef{耶 7:11} 所以，让我们选择属于生命的教会，好使我们得救。我想你们并非不知：活的教会是基督的身体（因为圣经说：\Quote{天主造了人，造了男造了女}\BibleRef{创 1:27}；男是基督，女是教会），而且圣经和宗徒都教导说，教会不属于现世，而是从起初就有的。因为它是属神的，正如我们的耶稣一样，但祂在末期才显现出来，为要拯救你们。\BibleRef{伯前 1:20} 教会既然是属神的，便在基督的肉身上显现，以此向我们表明：我们中若有人能在肉身中持守它而不败坏，他日后必在圣神中领受它。因为这肉身是那灵的预像；所以，凡败坏了预像的，日后必不能领受本物。因此，祂便这样说，弟兄们：\Quote{保存肉身，好使你们成为那灵的分享者。}若我们说肉身是教会、那灵是基督，那么，凡凌辱肉身的，便是对教会犯了凌辱之罪。这样的人，必不能分有那灵，即基督。这就是这肉身以后与圣神结合所能获得的永生和不朽；主为祂的选民所预备的，是人心不能述说、舌不能言表的。\BibleRef{格前 2:9}\par

% structure: p0030 source=h2 target=section reason=relative-heading-rank; base=h2

\section{第十五章  拯救者和被拯救者}

我不以为关于节制的劝告是无足轻重的；人若遵行，必不后悔，反能拯救自己和我这劝告者。\BibleRef{弟前 4:16}因为使一个迷失和将亡的灵魂回转，得以得救，不是小赏报。\BibleRef{雅 5:19}因为我们能将此种报偿献给创造我们的天主，即说话的和听道的都怀着信德和爱德说话和听道。因此，让我们持守我们因信成义成圣所行的道路，好能坦然无惧地向那说\Quote{你还在说话时，我必说：我在这里}的天主祈求。\BibleRef{依 58:9}因为这些话是大应许的记号，主说祂比求祂的人更乐意施予。我们既然分享了如此大的美善，就不要彼此嫉妒获得如此大的福分。因为这些话给听从者带来多大的喜乐，就给拒绝者带来多大的定罪。\par

% structure: p0032 source=h2 target=section reason=relative-heading-rank; base=h2

\section{第十六章 为审判之日做准备}

所以，弟兄们，我们既得了不小的悔改机会，就当趁有时机，转向那召叫我们的天主，趁着还有人接纳我们。因为如果我们弃绝这些放纵，克制灵魂，不满足它邪恶的欲望，我们就将分享耶稣的慈悲。要知道审判之日临近，如同燃烧的火炉，诸天和全地都要熔化，如铅在火中熔化；那时，人的隐密和显露的事都要显明。所以，施舍对于悔改罪过是好的；禁食胜于祈祷，施舍胜于二者；\Quote{爱德遮掩许多罪过，}\BibleRef{伯前 4:4}出于良心的祈祷救人脱免死亡。凡在这些事上齐全的人是有福的；因为施舍减轻罪的负担。\par

% structure: p0034 source=h2 target=section reason=relative-heading-rank; base=h2

\section{第十七章 续论同一主题}

因此，让我们全心悔改，免得我们中有人无谓地丧亡。因为如果我们领受了命令，并从事离弃偶像和教导他人的工作，那么一个已经认识天主的灵魂岂不更不该丧亡吗？所以，让我们互相帮助，扶助软弱的人行善，好使我们众人得救，彼此劝勉和告诫。不要只在受长老告诫时才相信和留意，当我们回到家中，也要记念主的诫命，不被世俗的欲望引诱回去，而要常常亲近，努力在主的命令上长进，好使我们同心合意，共同聚集得生命。因为主说：\Quote{我来聚集万民和异语的人。}这指的是祂显现的日子，那时祂要来救赎我们——各人按自己的行为受报。不信的人将看见祂的荣耀和大能，当他们看见耶稣拥有的世界权柄时，将惊讶地说：\Quote{我们有祸了，因为你存在，我们却不知道，不相信，也不听从那向我们明确指示我们救恩的长老。}\BibleRef{依 66:18}而\Quote{他们的虫子不死，他们的火也不灭；他们将成为一切血肉之躯所观看的景象。}\BibleRef{依 66:24}祂说的是大审判之日，那时他们将看见我们中间那些不敬虔、对耶稣基督的命令判断错误的人。义人既在忍受试炼和憎恶灵魂放纵上得胜，当他们看见那些偏离正路、言语或行为否认耶稣的人，在不可熄灭的火中受重刑惩罚时，将归光荣于自己的天主，说：\Quote{凡全心服事天主的，必有希望。}\par

% structure: p0036 source=h2 target=section reason=relative-heading-rank; base=h2

\section{第十八章 作者虽有罪，仍追求义}

因此，让我们也属于那感恩、服事天主的人，而不属于那受审判的不敬虔者。因为我自己虽全然是罪人，尚未逃离诱惑，仍处在魔鬼的工具中，却仍努力追求义，好能接近它，哪怕只是少许，因惧怕将来的审判。\par

% structure: p0038 source=h2 target=section reason=relative-heading-rank; base=h2

\section{第十九章 义人的赏报，即使他们受苦}

所以，弟兄姊妹们，我在真理的天主之后，向你们发出呼吁，要留意所写的话，好使你们拯救自己和那在你们中间宣读劝言的人。因为我所求的赏报是你们全心悔改，而你们则为自己获得救恩和生命。如此，我们便为一切愿意在敬虔和天主的良善上勤奋的年轻人树立榜样。我们不要因愚昧而不悦和恼怒，每当有人劝诫我们，使我们从不义转向义。因为有些恶行我们行了却不知道，是由于我们心中的三心二意和不信，我们的理智被虚妄的欲望弄瞎了。因此，让我们行义，好能得救到底。听从这些诫命的人是有福的，即使他们在今世暂时受苦，也将获得复活不朽的果实。所以，敬虔的人不要忧伤；即使目前受苦，那等候他的时间是有福的；复活与祖先同享生命，他将永远喜乐，毫无忧愁。\par

% structure: p0040 source=h2 target=section reason=relative-heading-rank; base=h2

\section{第二十章 敬虔，而非利得，才是真财富}

但你们不要因看见不义之人拥有财富，而天主的仆人却困窘，就心中不安。因此，弟兄姊妹们，让我们相信；我们在永生天主的考验中，在今世奋斗和受锻炼，好能在来世获得冠冕。没有义人立时得赏报，而是等候它。因为如果天主立即赏报义人，那么我们所行的就成了买卖，而非敬虔了。因为我们好像是追求利得而非敬虔而成为义人；因此，神圣的审判挫败了不义的心灵，并沉重地加上了锁链。\par

归于唯一的天主，不可见的，真理的圣父，祂曾派遣了救主和不朽的创始者给我们，借着祂也向我们显示了真理和天上的生命，愿光荣归于祂，直到永远。阿们。\par

\SourceRecord{Translated by John Keith.；Ante-Nicene Fathers；Vol. 9.；Edited by Allan Menzies.；Buffalo, NY: Christian Literature Publishing Co.,；1896.；<http://www.newadvent.org/fathers/1011.htm>.}
